% MIT OpenCourseWare: https://ocw.mit.edu
% 18.783 / 18.7831 Elliptic Curves Spring 2021
% License: Creative Commons BY-NC-SA 
% For information about citing these materials or our Terms of Use, visit: https://ocw.mit.edu/terms.
\newif\ifOCW
\OCWtrue
\documentclass[11pt]{article}
\usepackage{amsmath,amssymb,amsfonts}
\usepackage{listings}
\usepackage{courier}
\usepackage{enumerate}
\usepackage{hyperref}
\usepackage{color}
\definecolor{mylinkcolor}{rgb}{0.5,0.0,0.0}
\definecolor{myurlcolor}{rgb}{0.0,0.0,0.75}
\definecolor{mycitecolor}{rgb}{0.0,0.4,0.0}
\hypersetup{colorlinks=true,urlcolor=myurlcolor,citecolor=mycitecolor,linkcolor=mylinkcolor}
\ifOCW\usepackage{soul}\let\oldhref\href\renewcommand{\href}[2]{\oldhref{#1}{\ul{#2}}}\fi
\ifOCW\newcommand{\due}[1]{\hfill\phantom{Due: #1}}\else\newcommand{\due}[1]{\hfill{Due: #1}}\fi
\lstset{
	basicstyle=\small\ttfamily,
	keywordstyle=\color{blue},
	language=python,
	xleftmargin=16pt,
}

\textwidth=5.8in
\textheight=9in
\topmargin=-0.5in
\headheight=0in
\headsep=.5in
\hoffset=-.45in
\pagestyle{plain}

\newcommand{\Fp}{\mathbb{F}_p}
\newcommand{\Fq}{\mathbb{F}_q}
\newcommand{\FF}{\mathbb{F}}
\newcommand{\Q}{\mathbb{Q}}
\newcommand{\Z}{\mathbb{Z}}
\newcommand{\kron}[2]{\left(\frac{#1}{#2}\right)}
\newcommand{\SO}{\rm SO}
\newcommand{\SU}{\rm SU}
\newcommand{\M}{\textsf{M}}
\newcommand{\tr}{\operatorname{tr}}

\begin{document}
\setlength{\unitlength}{1in}
\begin{center}
\large
\textbf{18.783 Elliptic Curves\hspace{220pt}Spring~2021}\\\vspace{4pt}
\textbf{Problem Set \#2\due{03/05/2021}}\\\vspace{-6pt}
\normalsize
\begin{picture}(5.8,.1) 
\put(0,0) {\line(1,0){5.8}}
\end{picture}
\end{center}

\noindent\textbf{Description}: These problems are primarily related to material covered in Lecture 3.
\medskip

\noindent
\textbf{Instructions}: Solve any combination of problems summing to 100 points --- this will necessarily include the survey (Problem 6). Your solutions should be written up in latex (please do not submit handwritten solutions) and submitted as a pdf-file\ifOCW\else to \href{https://www.gradescope.com/courses/238945}{Gradescope}\fi.

Collaboration is permitted/encouraged, but you must identify your collaborators or your group\ifOCW\else\ on \href{https://psetpartners.mit.edu}{pset partners}\fi, as well any references you consulted that are not listed in the \ifOCW\href{https://ocw.mit.edu/courses/mathematics/18-783-elliptic-curves-spring-2021/syllabus}{syllabus}\else\href{http://math.mit.edu/classes/18.783/2021/syllabus.html}{syllabus}\fi\ or \ifOCW\href{https://ocw.mit.edu/courses/mathematics/18-783-elliptic-curves-spring-2021/lecture-notes-and-worksheets}{lecture notes}\else\href{http://math.mit.edu/classes/18.783/2021/lectures.html}{lecture notes}\fi. If there are none write ``\textbf{Sources consulted: none}'' at the top of your solutions.  Note that each student is expected to write their own solutions; it is fine to discuss the problems with others, but your writing must be your own.

The first person to spot each non-trivial typo/error in the problem set or lecture notes will receive 1-5 points of extra credit.
\smallskip

\noindent
\textbf{Note}: Several problems require a portion of your MIT ID as input.  If you would prefer not to use your MIT ID, let me know and I will choose a random 9-digit number that you can use in place of your MIT ID for the purpose of solving 18.783 problem sets.
\smallskip

\subsection*{Problem 1. Eight ways to Fibonacci (19 points)}
Let $F_k$ denote the $k$th Fibonacci number, defined by:
\begin{equation}\label{eq:fib}
F_1:=1,\qquad F_2:=1,\qquad F_{k}:=F_{k-2}+F_{k-1}\qquad(k\ge 2).
\end{equation}
Analyze the time complexity of computing $F_k$ modulo an $n$-bit prime~$p$ using each of the algorithms listed below.
You may assume that $k>n$.
Your answers should be in the form $O(f(k,n))$ and as tight as you can make them.
Use $\textsf{M}(n)$ to denote the complexity of multiplying two $n$-bit integers.
\begin{enumerate}[{\bf (a)}]
\item Compute the integers $F_1,F_2,\ldots, F_k$ via \eqref{eq:fib}, and then reduce $F_k$ mod $p$.
\item Same as (a), except compute in $\Fp$ rather than $\Z$ throughout.
\item Assume $p=\pm 1\bmod 5$.  Compute the roots $\phi$ and $\psi$ of $x^2-x-1$ in $\Fp$ (using a probabilistic root-finding algorithm), and then compute
$F_k=\frac{\phi^k-\psi^k}{\phi-\psi}$ in $\Fp$.
\item Assume $p\ne\pm 1\bmod 5$ and $p\ne 5$.  Represent $\FF_{p^2}$ as $\Fp[\phi]/(\phi^2-\phi-1)$, let $\psi=1-\phi$, and then compute $F_k=\frac{\phi^k-\psi^k}{\phi-\psi}$ as an element of $\FF_{p^2}$ that actually lies in $\Fp$.
\item Compute the $k$th power of $\begin{pmatrix}0&1\\1&1\end{pmatrix}$ in ${\rm GL}_2(\Fp)$ and output its upper right entry.
\item Assume $k$ is a power of 2.  Use the identities
\begin{align*}
F_{2i-1} &= F_i^2+F_{i-1}^2\\
F_{2i} &= (2F_{i-1}+F_i)F_i
\end{align*}
to compute $(F_1,F_2)$, $(F_3,F_4)$, $(F_7,F_8)$,\ldots $(F_{k-1},F_k)$, working modulo $p$ throughout.\footnote{This approach can be generalized to arbitrary $k$.}
\item
Notably absent from the above list is the standard recursive algorithm:
\begin{lstlisting}
def fibonacci(k):
    if k <= 2: return 1
    return fibonacci(k-1) + fibonacci(k-2)
\end{lstlisting}
What is the time complexity of computing \texttt{fibonacci(k)} and reducing it mod $p$?
What is the time complexity of computing \texttt{fibonacci(GF(p)(k))}?
\item Let $n:=10^6+N$, where $N$ is the last five digits of your student ID, and let $k=n^n$.\\  Compute $F_k\bmod 10007$.
\end{enumerate}


\subsection*{Problem 2. Cornacchia's algorithm (19 points)}
Cornacchia's algorithm computes primitive solutions $(x,y)$ to the Diophantine equation
\begin{equation}\label{eq:cornacchia}
x^2+dy^2 = m,
\end{equation}
where $d$ and $m$ are positive integers.  A \emph{primitive} solution has $x$ and $y$ relatively prime.
Typically $m=p$ or $m=4p$, where $p$ is a prime, but the algorithm works for any $m$, provided we are given an appropriate square root $r$ of $-d$ mod $m$ (if there are only two square roots, as when $m$ is prime $p$ or $4p$ with $p$ odd, it does not matter which we use, but in general one needs to check one of $\pm r$ for each square root $r$ of $-d$ modulo $m$).\footnote{Using the probabilistic root-finding algorithm described in \href{https://math.mit.edu/classes/18.783/2021/LectureNotes3.pdf}{Lecture 3} one can efficiently compute square roots modulo $p$, and with this one can also determine square roots modulo $4p$, and more generally, modulo any integer whose prime factorization is known.  But in general the problem of computing square roots modulo $m$ is believed to be as hard as factoring $m$.}



The algorithm uses a partial Euclidean algorithm that terminates as soon as the sequence of remainders $r_i$ drops below the square root of $r_0=m$.
\begin{enumerate}[1.]
\item
Let $r_0 = m$ and $r_1 = r$, where $r^2\equiv -d \bmod m$ and $0\le r\le m/2$.
\item
Compute $r_{i+2} = r_i \bmod r_{i+1}$ until $r_k^2 < m$ is reached.
\item
If $(m-r_k^2)/d$ is the square of an integer $s$, return the solution $(r_k,s)$.\\
Otherwise, return \texttt{null}.
\end{enumerate}

It is clear that if the algorithm returns $(r_k,s)$, then it is a solution to \eqref{eq:cornacchia}.
It is not so clear that the algorithm always finds a primitive solution if one exists, but this is true; see \cite{basilla} for a
short elementary proof.\footnote{There is an obvious typo in step 3 of the algorithm given in \cite{basilla} which is corrected above.}
If $m$ is square-free (as when $m$ is prime), every solution to \eqref{eq:cornacchia} is primitive, but this is not true in general (this is relevant to part (e)).

\medskip
\noindent
In this problem let $N:=n\cdot 10^{100}$, where $n$ is the last 4 digits of your student ID.

\begin{enumerate}[{\bf (a)}]
\item Implement this algorithm in Sage.  Use \texttt{mod(-d,m).is\_square()} to test if $-d$ has a square root mod $m$, and use \texttt{int(mod(-d,m).sqrt())} to get a square root.
\item You may recall Fermat's ``Christmas theorem", which states that an odd prime $p$ is the sum of two squares if and only if $p\equiv 1\bmod 4$.  You may also recall that $-1$ is a square modulo an odd prime $p$ if and only if $p\equiv 1\bmod 4$.

Let $n$ be the integer corresponding to the last 4 digits of your student ID.
For the least prime $p > N$ congruent to $1\bmod 4$, write $p$ as the sum of two squares.

\item Fermat also proved that a prime $p$ can be written in the form $p=x^2+3y^2$ if and only if $p\equiv 1\bmod 3$, which is equivalent to the condition that $-3$ is a square mod $p$.  For the least prime $p > N$ congruent to $1\bmod 3$, write $p$ in the form $p=x^2+3y^2$.

\item Show that this does not work for $d=5$ by finding a prime $p$ for which $-5$ is a square modulo $p$ but $p$ cannot be written in the form $x^2+5y^2$.  Empirically determine a stronger congruence condition on $p$ that guarantees not only that is $-5$ a square mod~$p$, but also that $p$ can be written in the form $x^2+5y^2$.
Then find the least prime $p >N$ that satisfies your condition and write $p$ in the form $p=x^2+5y^2$.

\item Let $E$ be the elliptic curve $y^2=x^3-35x-98$ and let $p\ne 2,7$ be a prime.  Like the elliptic curves considered in Problem 5 of Problem Set 1, the elliptic curve $E$ has complex multiplication, and the integer $a_p=p+1-\#E(\Fp)$ is zero if and only if $-7$ is not a square modulo $p$.
When $-7$ is a square modulo $p$, the integer $a_p$ satisfies the equation $4p=a_p^2+7y^2$, for some $y\in\Z$.
Prove that this equation has a solution $(a_p,y)$ if and only if the equation $p=u^2+7v^2$ has a solution $(u,v)$.

Use your algorithm to find a solution to $p=u^2+7v^2$ for the least prime $p>N$ for which $-7$ is a square modulo $p$,
and use this to deduce the absolute value of $a_p$.
Determine the sign of $a_p$, by finding a random point $P\in E(\Fp)$ for which only one of $(p+1-a_p)P$ and $(p+1+a_p)P$ is zero.

\end{enumerate}

\subsection*{Problem 3. Computing $r$th roots in cyclic groups (38 points)}
In Lecture 3 we saw how to compute $r$th roots in a finite field $\Fq$ using a probabilistic root-finding algorithm.
In this problem you will implement an entirely different approach for computing $r$th roots that works in any finite cyclic group $G$, including $G=\Fq^\times$.
In addition to being more general, this method is typically faster than using probabilistic root-finding to compute $r$th roots in $\Fq^\times$ (but this depends on the values of $r$ and  $q$).

We assume without loss of generality that $r$ is prime (to compute $n$th roots, successively compute $r$th roots for the primes $r$ dividing $n$, with multiplicity).
To simplify notation we write $G$ additively, so an \emph{$r$th root} of $\gamma\in G$ is an element $\rho\in G$ for which $r\rho=\gamma$.
Let $|\gamma|$ denote the \emph{order} of~$\gamma$, the least positive integer~$m$ for which $m\gamma=0$.

\medskip
\noindent
\textbf{Prove the following statements}:
\begin{enumerate}[{\bf (a)}]
\item For all $\gamma\in G$ and $n\in\Z$ we have $|n\gamma| = |\gamma| / \gcd(n,|\gamma|)$.
\item If $r$ does not divide $|G|$, then there is an integer $s$ such that for all $\gamma\in G$ the element $\rho = s\gamma$ is the unique $r$th root of $\gamma$.
\item  If $r$ does divide $|G|$, then the number of $r$th roots of each $\gamma\in G$ is either 0 or $r$.
In the latter case, the $r$th roots of $\gamma$ do not necessarily lie in $\langle \gamma\rangle$ (give an example).
\item Suppose $r$ divides $|G|$.  Let $|G|=ar^k$, where $r\nmid a$.  Let $\delta\in G$ be an element of order~$r^k$, let $\gamma$ be any element of $G$, and $\alpha=a\gamma$ and $\beta=r^k\gamma$.  The following hold:

\begin{enumerate}[(i)]
\setlength{\itemsep}{0pt}
\item $\alpha = x\delta$ for some integer $x\in[1,r^k]$.
\item If $r$ does not divide $x$ then there is no $\rho\in G$ for which $r\rho = \gamma$.
\item If $r$ divides $x$, and $s$ and $t$ are integers satisfying $sa+tr^{k+1}=1$,
then the element $\rho = s(x/r)\delta + t\beta$ satisfies $r\rho = \gamma$.
\end{enumerate}
\end{enumerate}
\vspace{-4pt}

\noindent
The element $\delta$ is a generator for the $r$-Sylow subgroup of~$G$.  If $G=\langle\sigma\rangle$, we can use $\delta=a\sigma$.
The integer $x$ is the \emph{discrete logarithm} of $\alpha$ with respect to $\delta$.

\bigskip
\noindent
\textbf{Implement the following algorithm} for computing an $r$th root of $\gamma$ in a cyclic group~$G$ of order $ar^k$, where $r$ is a prime that does not divide $a$, given $\delta\in G$ of order $r^k$:
\begin{enumerate}
\setlength{\itemsep}{-1pt}
\item If $k=0$ then compute $s=1/r\bmod a$ and return $\rho = s\gamma$.
\item Compute $\alpha=a\gamma$ and $\beta=r^k\gamma$.
\item Compute the discrete logarithm $x$ of $\alpha$ with respect to $\delta$ by brute force: check whether $\alpha=x\delta$ for each $x$ from 1 to $r^k$ (this holds for some $x$, by part (a) of 4).\footnote{We will learn much better ways to compute this discrete logarithm later in the course.  For the moment, assume $r^k$ is small (for finite fields $\Fq$, this is usually the case, even when $q$ is very large).}
\item If $r$ does not divide $x$ then return \texttt{null}.
\item Compute $s$ and $t$ such that $sa+tr^{k+1}=1$ using the extended Euclidean algorithm.
\item Return $\rho=s(x/r)\delta+t\beta$.
\end{enumerate}
The return value \texttt{null} is used to indicate that $\gamma$ does not have any $r$th roots in $G$.
To compute $s=1/r\bmod a$ in Sage, use: \texttt{s=1/mod(r,a)}.
To compute $s$ and~$t$ such that $sa+tr^{k+1}=1$, use: \texttt{d,s,t=xgcd(a,r**(k+1))} (here $d=\gcd(a,r^{k+1})$ is 1).

The Python language used by Sage is untyped, so your algorithm can be used to compute $r$th roots in any cyclic group that Sage knows how to represent; it will automatically perform operations in whatever group the inputs $\delta$ and $\gamma$ happen to lie in.
To test your algorithm, you may find it useful to work in the additive group of the ring $\Z/n\Z$, where $n=ar^k$, which you can create in Sage using \texttt{Zn=Integers(n)}.
You can then use \texttt{delta=Zn(a)} to create an element of $\Z/n\Z$ with additive order $r^k$.

\bigskip

Let $E$ be the elliptic curve $y^2=x^3+31415926x+27182818$ over $\Fp$ with $p=2^{255}-19$.
The group $E(\Fp)$ is cyclic, of order $n = 2\cdot 3\cdot 31\cdot m$, where
\small
$$m=311269057089559665117126303786795451217418463436862985689835777395934466489,$$
\normalsize
and the point $P=(x_0:y_0:1)$, where $x_0=99$ and
\small
$$y_0=3646051633135286488902046129458077014725501801396015176760137375427748642285,$$
\normalsize
is a generator for $E(\Fp)$.  Thus for $r=2,3,31$ you can use $\delta = (n/r)P$ as a generator of the $r$-Sylow subgroup (which in each case has order $r$).

Let $c$ be the least prime greater than the integer formed by the last four digits of your student ID.
Let $Q=(x_1:y_1:1)$, where
\small
\begin{align*}
x_1&=43125933575059134974422288266359854378815207690220011740187158431378585841262,\\
y_1&=30438392960540783858586956956150489842875282144799753811252714114065692010946
\end{align*}
\normalsize
\begin{enumerate}
\item[{\bf (e)}]
Use your algorithm to find an $r$th root $R$ of $\gamma=cQ$, for $r=2,3,31$.
Note that you can easily check your result by testing whether \texttt{r*R==c*Q} holds using Sage (please be sure to do this).
In your answer you only need to list the point $R$ for each value of $r$, you don't need to include your code.
Be sure to format your answer so that the coordinates of $R$ all fit on the page.
\end{enumerate}

\subsection*{Problem 4. Exponentiating with addition chains (38 points)}
An \emph{addition chain} for a positive integer $n$ is an increasing sequence of integers $(c_0,\ldots,c_m)$ with $c_0=1$ and $c_m=n$ such that each entry other than $c_0$ is the sum of two (not necessarily distinct) preceding entries.
The \emph{length} of an addition chain is the index $m$ of the last entry.
When computing $a^n$ with a generic algorithm, the exponents of the powers of $a^k$ computed by the algorithm define an addition chain whose length is the number of multiplications performed.
For example, using left-to-right binary exponentiation to compute $a^{47}$ yields the addition chain $(1,2,4,5,10,11,22,23,46,47)$,
and right-to-left binary exponentiation yields the addition chain $(1,2,3,4,7,8,15,16,32,47)$, both of which have length~9.

\begin{enumerate}[{\bf (a)}]
\item For $n=715$, determine the addition chains given by: (i) left-to-right binary, (ii) right-to-left binary, (iii) fixed-window, and (iv) sliding-window exponentiation, using a window of size 2 for (iii) and~(iv).
\item Find an addition chain for $n=715$ that is shorter than any you found in part (a).
\item Repeat part (a) for the integer $N$ obtained by adding 990,000 to the last 4 digits of your student ID, using a window of size 3 for the fixed and sliding window cases.
\item Find the shortest addition chain for $N$ that you can.
There is a good chance you can do better than any of the chains you found in part (c).
\end{enumerate}

In groups where inversions are cheap (such as the group of points on an elliptic curve), it is advantageous to use \emph{signed} binary representations of exponents, where we write the exponent $n$ in the form $n=\sum n_i2^i$ with $n_i\in\{-1,0,1\}$.  Such a representation is generally not unique, but there is a unique signed representation with the property that no pair of adjacent digits are both nonzero.  This is known as \emph{non-adjacent form} (NAF).  The NAF representation of 47, for example, is $10\bar{1}000\bar{1}$, where $\bar{1}$ denotes $-1$.

To construct the NAF representation one begins by writing $n$ in binary with a leading~$0$, and then successively replaces the least significant block of the form $01\cdots1$ with $10\cdots0\bar{1}$ until there are no adjacent nonzero digits.  For example, the computation for $47$ proceeds as $0101111, 11000\bar{1}, 10\bar{1}000\bar{1}$, which reduces the number of nonzero digits from 5 to 3.  Even though the length is increased by 1, the total cost of exponentiation may be reduced.

An \emph{addition-subtraction} chain extends the definition of an addition chain by allowing $c_k=c_i\pm c_j$.
Exponentiation using the NAF representation defines an addition-subtraction chain.
For example, using left-to-right binary exponentiation, the NAF representation of 47 yields the chain $(1,2,4,3,6,12,24,48,47)$,
which is shorter than the addition chain $(1,2,4,5,10,11,22,23,46,47)$ given by standard left-to-right binary exponentiation.

\begin{enumerate}[{\bf (a)}]
\setcounter{enumi}{4}
\item Compute addition-subtraction chains for $n=715$ and the integer $N$ defined in part (c) using left-to-right binary exponentiation with the NAF representation.
\item Find the shortest addition-subtraction chains for $n$ and $N$ that you can.
\end{enumerate}

\subsection*{Problem 5. Root-finding over $\Z$ (76 points)}
In this problem you will develop an algorithm to find integer roots of polynomials in $\Z[x]$ using a $p$-adic version of Newton's method (also known as Hensel lifting).
As an application, this gives us an efficient way to factor perfect powers (a special case that we will need to handle when we come to the elliptic curve factorization method),
and it will be used as a black box in Problem 2 to find integer roots of division polynomials.

In the questions below, $p$ can be any integer greater than 1, but you may assume it is a prime power if you wish.

\begin{enumerate}[{\bf (a)}]
\item Let $x_0\in\Z$ and $f\in\Z[x]$.  Prove that the following equivalence holds in $\Z[x]$:
\[
f(x) \equiv f(x_0) + f'(x_0)(x-x_0)\bmod (x-x_0)^2.
\]
\item
Let $x_0,z_0\in\Z$ and $f\in\Z[x]$ satisfy $f(x_0)\equiv 0\bmod p$ and $f'(x_0)z_0\equiv 1\bmod p$.
Let
\begin{align*}
x_1 &\equiv x_0-f(x_0)z_0 \bmod p^2,\\
z_1 &\equiv 2z_0-f'(x_1)z_0^2\bmod p^2.
\end{align*}
Prove that that the following three equivalences hold:
\begin{align}
x_1&\equiv x_0\bmod p,\tag{i}\\
f(x_1)&\equiv 0\bmod p^2,\tag{ii}\\
f'(x_1)z_1&\equiv 1\bmod p^2.\tag{iii}
\end{align}
Show that (i) and (ii) characterize $x_1\bmod p^2$ uniquely by proving that if $x_2\in\Z$ also satisfies $x_2\equiv x_0\bmod p$ and $f(x_2)\equiv 0\bmod p^2$, then $x_1\equiv x_2\bmod p^2$.
\end{enumerate}

\noindent
Iteratively applying (b) yields an algorithm that, given an integer $k$ and $x_0,z_0$, and~$f$ satisfying the hypothesis (a), outputs an integer $x_k$ that satisfies $f(x_k)\equiv 0\bmod p^{2^k}$.

\begin{enumerate}[{\bf (a)}]
\setcounter{enumi}{2}
\item
Prove that if $f$ has an integer root $r$ for which $f'(r)$ is invertible modulo $p$, then given $x_0\equiv r\bmod p$, $z_0\equiv 1/f'(x_0)\bmod p$, and $k$ such that $|r| < p^{2^k}/2$, this algorithm outputs $x_k$ such that $r$ is the unique integer $r\equiv x_k \bmod p^{2^k}$ satisfying $|r| < p^{2^k}/2$.
\end{enumerate}

\noindent
To apply the result in (c), we need to know a suitable starting value (or values) for $x_0$.  For the two applications we have in mind, this is will be straightforward, so let us proceed on the assumption that we are given a suitable $x_0$ and $z_0$ such that $f'(x_0)z_0\equiv 1\bmod p$.

Let $B$ be the maximum of the absolute values of the coefficients of~$f$, and let $B_0$ be an upper bound on the absolute value of its largest integer root.  It suffices to choose the least~$k$ such that $p^{2^k} > 2B_0$, and since any integer root of $f$ must divide its constant coefficient, we can assume that $B_0\le B$.
We can also assume $p < 2B_0$, since otherwise the problem is trivial ($k=0$ and $x_k=x_0$).

\begin{enumerate}[{\bf (a)}]
\setcounter{enumi}{3}
\item
Prove that with this choice of $k$ the algorithm can be implemented to run in time $O(d\ \M(\log B))$, where $d$ is the degree of~$f$ (be careful here, the most obvious implementation will not achieve this time bound).
Prove that if $f$ has $O(1)$ terms, then the algorithm can be implemented to run in time $O(\M(\log B) + \M(\log B_0)\log d)$.

\item
Using the primes $p=2$ and $p=3$, describe an efficient algorithm that, given an integer $N$ relatively prime to 6, either outputs an integer~$a$ and a prime~$q$ such that $a^q=N$, or proves that $N$ is not a perfect power.
Prove that your algorithm runs in quasi-quadratic time (meaning $O(n^2(\log n)^c)$ for some constant $c$, where $n=\log N$).
\end{enumerate}

\begin{enumerate}[{\bf (a)}]
\setcounter{enumi}{5}
\item
Implement your algorithm and report the result and running time on the each of the following inputs:
$2^{1000}+297,\ 5^{503},\ (2^{500}+55)^2,\ (2^{333}+285)^3$, $(2^{32}+15)^{31}$, $500!$.
To time your code in sage, use the \texttt{timeit} function (e.g. \texttt{timeit("1+1")}), or in a CoCalc Jupyter notebook, use the \texttt{\%timeit} \href{https://share.cocalc.com/share/8b892baf91f98d0cf6172b872c8ad6694d0f7204/PythonDataScienceHandbook/notebooks/01.07-Timing-and-Profiling.ipynb}{line-magic}.
\item
Prove that the algorithm you gave in (e) can be implemented to run in sub-quadratic time (meaning $o(n^2)$ where $n=\log N$).  You may need to modify your algorithm in order to achieve this.\footnote{In fact, this problem can be solved in quasi-linear time \cite{bernstein}.}
\end{enumerate}

\subsection*{Problem 6. Survey (5 points)}
Complete the following survey by rating each of the problems you attempted on a scale of 1 to~10 according to how interesting you found the problem (1 = ``mind-numbing," 10 = ``mind-blowing"), and how difficult you found it (1 = ``trivial," 10 = ``brutal").  Also estimate the amount of time you spent on each problem to the nearest half hour.

\begin{center}
\begin{tabular}{l|r|r|r|}
& Interest & Difficulty & Time Spent\\\hline
Problem 1 & & & \\\hline
Problem 2 & & & \\\hline
Problem 3 & & & \\\hline
Problem 4 & & & \\\hline
Problem 5 & & & \\\hline
\end{tabular}
\end{center}
\noindent
Also, please rate each of the following lectures that you attended, according to the quality of the material (1=``useless", 10=``fascinating"), the presentation (1=``epic fail", 10=``perfection"), the pace (1=``way too slow", 10=``way to fast"), and the novelty of the material (1=``old hat", 10=``all new").

\begin{center}
\begin{tabular}{l|l|r|r|r|r|}
Date & Lecture Topic & Material & Presentation & Pace & Novelty\\\hline
2/24 & Finite field arithmetic & & & & \\\hline 
3/1 & Isogenies & & & & \\\hline 
\end{tabular}
\end{center}

\noindent
Feel free to record any additional comments you have on the problem sets or lectures.

\begin{thebibliography}{9}
\bibitem{basilla}
Julius Magalona Basilla, \href{http://projecteuclid.org/euclid.pja/1116442240}{\textit{On the solution of $x^2+dy^2=m$}}, Proc. Japan Acad. Ser. A Math Sci. \textbf{80} (2004), 40--41.
\bibitem{bernstein} D. J. Bernstein, \href{http://www.ams.org/journals/mcom/1998-67-223/S0025-5718-98-00952-1/}{\textit{Detecting perfect powers in essentially linear time}}, Math. Comp. \textbf{67} (1998), 1252--1283.
\end{thebibliography}

\end{document}